\chapter{地平線、状態方程式、真空エネルギー}

\section{優先座標系、超光速後退、地平線}

宇宙マイクロ波背景放射は、確かに特別な座標系を選び出す。その放射が等方的に見える座標系では、近傍の銀河もまた平均として静止している。我々が背景放射に対して運動すると、前方の光子はドップラー偏移を受けてより高いエネルギーになり、後方の光子はより低いエネルギーになる。優先座標系とは、マイクロ波背景放射が等方的に見える座標系であり、それはまた平均として、局所的な銀河が静止している座標系でもある。

\begin{classroomqa}
\lecturetimestamp{00:00:25}
\audiencequestion{宇宙放射の見かけの波長は観測者の運動に依存するのか、そしてそれは優先座標系の存在を意味するのか。}
\lecturerresponse{はい。マイクロ波背景放射を通り抜ける運動はドップラー異方性を生み出すので、優先的な宇宙論的座標系が存在する。すなわち、マイクロ波背景放射が等方的に見える座標系である。それはまた平均として、局所的な銀河が静止している座標系でもある。}
\end{classroomqa}

第二の問いは後退速度に関する。ハッブルの法則によれば、十分に遠い銀河は次に従って互いに離れていく。
\begin{equation}
v = H d .
\end{equation}
十分に大きな\(d\)に対しては、この後退速度は光速を超える。これは相対性理論と矛盾しない。なぜなら禁止は局所的なものであり、一つの物体が我々のすぐ目の前で別の物体を光より速く追い越すことはできないからである。遠方の銀河についての言明は、広く離れた場所の間の幾何学についての言明である。

\begin{classroomqa}
\lecturetimestamp{00:04:48}
\audiencequestion{二つの銀河は\(c\)より大きな相対速度を持ちうるか。}
\lecturerresponse{はい、それらが十分に離れていれば可能である。ハッブル後退は局所的な相対性理論に違反することなく光速を超えうるし、後退速度が\(c\)に達する場所が地平線である。}
\end{classroomqa}

後退が光速に達すると、我々に向けて送られた信号はもはや正味の前進をしなくなる。光線は依然として局所的に限界速度で進むが、克服しなければならない距離があまりに速く増大するのである。

\begin{classroomqa}
\lecturetimestamp{00:07:03}
\audiencequestion{逆方向に送られた光線なら、依然として我々に向かって進むのではないか。}
\lecturerresponse{局所的にはそうである。しかし空間は光線が克服できるよりも速く膨張する。光線は近傍の点に対して限界速度で進むが、我々へ戻る全距離は増大し続ける。}
\end{classroomqa}

\section{膨張する計量と線の類推}

有用な描像は一次元のものである。原子からなる線を想像しよう。普通の物質の線を引き伸ばせば、やがてそれは切れる。ここで描像を修正する。隣り合う原子が離れすぎるたびに、我々はそれらの間に別の原子を挿入する。線は成長し続けるが、その局所的な性質は薄まらない。

\begin{figure}
\centering
\includegraphics[width=0.60\textwidth]{lecture_04_figure_02.png}
\caption{膨張する計量を描くために用いられる一次元の線。}
\end{figure}

この描像においても、依然として局所的な限界信号速度が存在する。ちょうど普通の線の上の音が限界速度を持つように。規則は局所的である。すなわち、撹乱は近傍の点に対してその速度を追い越すことはできない。それにもかかわらず、線が十分に長く、新しい点が十分に速く挿入されれば、その十分に遠い部分は信号が決して戻れないほど速く後退する。

したがって、特殊相対論的な速度の合成は話の一部にすぎない。それは観測者が同じ点ですれ違うとき、局所的に適用される。膨張する幾何学における広く離れた銀河に対しては、正しい道具ではない。

\begin{classroomqa}
\lecturetimestamp{00:09:42}
\audiencequestion{この状況でも、速度を合成する通常の相対論的法則は依然として成り立つのか。}
\lecturerresponse{局所的には、ある点においては、はい。観測者が次々とすれ違うなら、通常の特殊相対論的規則が適用される。しかしそれは、\(c\)を超える広く離れた後退を記述する正しい方法ではない。そこでは膨張する空間の幾何学に立ち返らなければならない。}
\end{classroomqa}

同じ描像が赤方偏移の簡単な説明を与える。線の上に波を描く。線が膨張するにつれて、波長もそれとともに膨張する。膨張する幾何学において、これは単に示唆的なイメージではない。膨張する時空における波動方程式は、実際に引き伸ばされた波長を生み出すのである。

\begin{classroomqa}
\lecturetimestamp{00:14:42}
\audiencequestion{波が宇宙とともに引き伸ばされるという言明は、理論的なものか、経験的なものか、それとも両方か。}
\lecturerresponse{それは膨張する時空における波の伝播の理論的帰結であり、また天文学的観測によっても支持されている。}
\end{classroomqa}

この描像はゴムバンドではない。ゴムバンドは引き伸ばされると薄くなり、そのような媒質に住む者は局所的な性質が変化するのを見るであろう。ここでは、大域的な距離の尺度が成長する一方で、局所的な性質は同じままであり続ける。

\section{最終散乱と観測可能な宇宙}

地平線とは、後退速度が光速に達する場所である。最終散乱面はそのわずかに内側にある。地平線から放出された信号は無限に引き伸ばされるであろうが、最終散乱放射の赤方偏移は巨大ではあるが有限である。

\begin{classroomqa}
\lecturetimestamp{00:17:53}
\audiencequestion{超光速後退と最終散乱面と観測可能な宇宙の関係は何か。}
\lecturerresponse{地平線とは、後退速度が光速に達する場所である。最終散乱面はそのわずかに内側にある。その赤方偏移が巨大ではあるが有限だからである。実用上、それは観測可能な宇宙の端を画す。}
\end{classroomqa}

最終散乱放射の赤方偏移はおよそ\(10^3\)倍である。これはすでに非常に大きいので、最終散乱面は多くの目的にとって地平線に極めて近い。観測可能な宇宙とは、光が我々に到達しうる宇宙の部分である。

後期には、膨張は真空優勢の領域に近づく。するとハッブル量は定数に近づき、地平線は一定の大きさに近づき、銀河もマイクロ波放射も次第に視界から失われていく。

\begin{classroomqa}
\lecturetimestamp{00:22:43}
\audiencequestion{\(H\)は時間とともに変化するので、地平線距離は変化しているのか、そしてマイクロ波背景放射は最終的に消えるのか。}
\lecturerresponse{後期には\(H\)は定数に近づくので、地平線は一定の距離に近づく。マイクロ波背景放射は引き伸ばされ、冷え、そして薄れ続ける。遠方の銀河もまた地平線を通過し、視界から消える。}
\end{classroomqa}

すでに重力で束縛されている構造だけが可視のまま残る。太陽系はばらばらに飛び散らず、銀河もばらばらに飛び散らないが、大規模な宇宙は孤立した島となる。

\section{基本膨張方程式と三つの成分}

ここで基本膨張方程式は
\begin{equation}
\left(\frac{\dot a}{a}\right)^2 = \frac{8\pi G}{3}\rho ,
\qquad c=1 .
\end{equation}
である。ここで\(a(t)\)はスケール因子であり、\(\rho\)は全エネルギー密度である。

\begin{figure}
\centering
\includegraphics[width=0.74\textwidth]{lecture_04_figure_03.png}
\caption{基本膨張方程式と一次元の線の類推。}
\end{figure}

力学的に重要な三つの成分がある。

物質については、
\begin{equation}
\rho=\frac{\text{const.}}{a^3},
\qquad
a(t)\sim t^{2/3}.
\end{equation}
これは、全体の膨張とともに運動し、共動系では本質的に静止している質量をもつ物体の場合である。

放射については、
\begin{equation}
\rho=\frac{\text{const.}}{a^4},
\qquad
a(t)\sim t^{1/2}.
\end{equation}
余分な\(a\)の冪は、宇宙が膨張するにつれて各光子が赤方偏移することに由来する。

また、
\begin{equation}
\rho_0=\text{constant},
\end{equation}
真空エネルギーがあり、暗黒エネルギーとも呼ばれ、慣用的な数因子を除けば宇宙定数でもある。この成分はまったく希薄化しない。

次の段階は、エネルギー密度を圧力に関係づけることである。

\section{放射圧と冷たい物質}

放射から始めよう。長さ\(L\)の線分にすぎない一次元の箱を考え、その中で光子が往復して跳ねているとする。一次元では、力と圧力は事実上同じものである。

運動量の大きさ\(p\)の光子をとる。それが壁で反射するとき、その運動量は
\begin{equation}
\Delta p = 2p.
\end{equation}
だけ変化する。同じ壁との連続する衝突の間の時間は
\begin{equation}
\Delta t = \frac{2L}{c}.
\end{equation}
である。したがって壁に働く平均の力は
\begin{equation}
F=\frac{\Delta p}{\Delta t}=\frac{pc}{L}.
\end{equation}
である。光子については
\begin{equation}
E=pc,
\end{equation}
なので
\begin{equation}
F=\frac{E}{L}.
\end{equation}
となる。エネルギーを長さで割ったものは一次元のエネルギー密度である。この一次元の場合、
\begin{equation}
P=\rho .
\end{equation}

多数の光子があっても、平均をとれば同じ関係が生き残る。全力は単位長さあたりの全エネルギーにほかならない。

次に三次元に移ろう。もしすべての光子が一つの軸に沿ってのみ運動するなら、それらは一対の面に圧力を及ぼし、他の面には及ぼさないであろう。これは故意に誤っているが、何を修正しなければならないかを直ちに示している。等方的な光子の海では、運動量はすべての方向に分布している。平均して、エネルギー密度の三分の一が各圧力成分に寄与する。結果は
\begin{equation}
P=\frac{\rho}{3}.
\end{equation}
である。

冷たい物質については答えは異なる。共動する箱、すなわち宇宙の流れとともに運動する小さな領域をとる。内部の粒子はエネルギーを運んでいるが、もしそれらが箱に対して静止していれば、箱の壁にぶつかることはない。したがって圧力は
\begin{equation}
P=0.
\end{equation}
である。

\begin{figure}
\centering
\includegraphics[width=0.74\textwidth]{lecture_04_figure_04.png}
\caption{物質と放射：\(a\sim t^{2/3}\)、\(P=0\)、および\(a\sim t^{1/2}\)、\(P=\rho/3\)。}
\end{figure}

板書のまとめは
\begin{align}
\text{matter:}\qquad & a(t)\sim t^{2/3}, \qquad P=0,\\
\text{radiation:}\qquad & a(t)\sim t^{1/2}, \qquad P=\frac{\rho}{3}.
\end{align}
である。

\begin{classroomqa}
\lecturetimestamp{00:54:27}
\audiencequestion{膨張する宇宙では、粒子は実際には互いに離れているのではないか。}
\lecturerresponse{そのとおりだが、圧力の議論は無限小の共動する箱に適用される。その局所的なスケールでは、熱力学的な推論は依然として成り立つ。}
\end{classroomqa}

\section{状態方程式}

圧力はエネルギー密度に比例するとみなす。
\begin{equation}
P=w\rho .
\end{equation}
黒板では同じパラメータが\(\omega\)と書かれている。二つの基本的な場合については、
\begin{equation}
w=0 \quad \text{for cold matter},
\qquad
w=\frac{1}{3} \quad \text{for radiation}.
\end{equation}

中間的な場合もあり得る。\(0\)と\(1/3\)の間の値は、完全に冷たいわけでも完全に相対論的でもない物質を表す。ここで挙げられた例は温かい暗黒物質である。

\begin{figure}
\centering
\includegraphics[width=0.74\textwidth]{lecture_04_figure_05.png}
\caption{状態方程式\(P=\omega\rho\)と密度のスケーリングの議論。}
\end{figure}

次の問題は、\(w\)が\(a\)に対する\(\rho\)の依存性を支配するかどうかである。答えは熱力学的な恒等式から得られる。

体積\(V\)の箱を考える。中身が圧力\(P\)を及ぼすなら、箱が\(dV\)だけ膨張するとき、壁に対してなされる仕事は
\begin{equation}
dE=-P\,dV.
\end{equation}
ここでエネルギーを
\begin{equation}
E=\rho V.
\end{equation}
と書く。その微分は
\begin{equation}
dE=\rho\,dV+V\,d\rho .
\end{equation}
二つの表式を等しいと置くと
\begin{equation}
\rho\,dV+V\,d\rho=-P\,dV,
\end{equation}
すなわち
\begin{equation}
V\,d\rho=-(\rho+P)\,dV.
\end{equation}
\(P=w\rho\)を用いると、
\begin{equation}
V\,d\rho=-(1+w)\rho\,dV.
\end{equation}
\(\rho V\)で割ると、
\begin{equation}
\frac{d\rho}{\rho}=-(1+w)\frac{dV}{V}.
\end{equation}
積分すると、
\begin{equation}
\rho \propto V^{-(1+w)}.
\end{equation}
共動宇宙論的箱に対しては、
\begin{equation}
V\propto a^3,
\end{equation}
したがって
\begin{equation}
\rho \propto a^{-3(1+w)}.
\end{equation}

先の場合が直ちに再現される。
\begin{align}
w=0 &\Longrightarrow \rho \propto a^{-3},\\
w=\frac{1}{3} &\Longrightarrow \rho \propto a^{-4}.
\end{align}

これが\(w\)の重要な理由である。\(\rho(a)\)が分かれば、それをフリードマン方程式に戻すことができる。次に冪則\(a\sim t^\alpha\)を試し、その状態方程式に対して宇宙がどのように膨張するかを決定する。

\section{真空エネルギーと負の圧力}

一つの特別な場合は異なる振る舞いをする。
\begin{equation}
w=-1.
\end{equation}
すると
\begin{equation}
1+w=0,
\end{equation}
なので指数は消え、
\begin{equation}
\rho=\rho_0=\text{constant}.
\end{equation}
対応する圧力は
\begin{equation}
P=-\rho .
\end{equation}

負の圧力それ自体は不思議ではない。正の圧力は箱の壁を外へ押す。負の圧力は内へ引く。それは張力である。引き伸ばしたゴムバンドが日常的な例を与える。それは正のエネルギーを蓄えるが、両端を引き寄せる。

一定密度\(\rho_0\)をフリードマン方程式に代入する。
\begin{equation}
\left(\frac{\dot a}{a}\right)^2=\frac{8\pi G}{3}\rho_0.
\end{equation}
右辺は単なる数なので、
\begin{equation}
\frac{\dot a}{a}=\sqrt{\frac{8\pi G}{3}\rho_0}.
\end{equation}
これはハッブル量が本当に一定である特別な場合である。方程式は
\begin{equation}
\dot a=\sqrt{\frac{8\pi G}{3}\rho_0}\,a,
\end{equation}
と書くことができ、解は
\begin{equation}
a(t)\propto \exp\!\left(\sqrt{\frac{8\pi G}{3}\rho_0}\,t\right).
\end{equation}

これは指数関数的膨張である。関係する時間スケールは倍加時間ではなく\(e\)倍時間であり、ここで用いられた粗い見積もりではそれを宇宙の年齢程度に置く。真空エネルギーと宇宙定数は、歴史的な数値の慣習を別にすれば同じ物理的要素である。

\begin{classroomqa}
\lecturetimestamp{01:27:16}
\audiencequestion{真空エネルギーが指数関数的膨張を駆動するなら、やがてすべてが引き裂かれてしまうのではないか。}
\lecturerresponse{必ずしもそうではない。束縛系はそれに抵抗できる。銀河、太陽系、あるいは原子は引き裂かれる必要はないが、十分に大きく希薄な構造は膨張に参加し、やがて離れていく。}
\end{classroomqa}

\begin{classroomqa}
\lecturetimestamp{01:31:43}
\audiencequestion{フリードマン方程式の平方根を取るとき、負の符号に何か興味深いものがあるか。}
\lecturerresponse{はい。平方根は二つの分枝を与える。一つは膨張解であり、もう一つは時間反転した収縮解である。方程式自体は両方を許す。}
\end{classroomqa}

\section{三つの時代と観測された膨張史}

三つの成分がすべて存在し、それらは異なるスケーリングをする。
\begin{equation}
\begin{aligned}
\rho_{\text{radiation}}&\propto a^{-4},\\
\rho_{\text{matter}}&\propto a^{-3},\\
\rho_0&=\text{constant}.
\end{aligned}
\end{equation}
これは直ちに三つの時代を意味する。

非常に小さい\(a\)では、放射が支配する。その後、物質が支配する。十分に大きい\(a\)では、真空エネルギーが支配する。対応する膨張則は
\begin{equation}
a(t)\sim t^{1/2},
\qquad
a(t)\sim t^{2/3},
\qquad
a(t)\sim e^{Ht}.
\end{equation}

ここで挙げられた粗い時間は次のとおりである。放射は約\(10^4\)年の後に物質に道を譲る。宇宙はその後、約\(10^5\)年の後に透明になる。物質支配から加速膨張への転移はさらにずっと後、数十億年の後に起こり、用いられている桁の見積もりでは宇宙史のほぼ中間である。

この歴史はどのように推測されるのか。我々は未来を見ることはできないが、過去を見ることはできる。より大きな距離を見ることは、より早い時間を見ることである。遠方の銀河を観測し、その運動が距離とどう相関するかを測定し、ハッブル量を後方視線時間の関数として再構成することによって、スケール因子の時間履歴を再構成する。放射支配の時代は、通常の銀河天文学で直接見るには早すぎるが、物質支配の時代と真空支配への転移は、再構成された曲線に見ることができる。

\section{遠い未来における最後の散乱}

未来のマイクロ波の空は、銀河が地平線の向こうへ消えていく話とは同じではない。より後の時間に見る放射は、より遠方の物質から来るので、効果は同じ原子が単に遠ざかるのを見ることではなく、より深くより冷たいものを見ることである。

\begin{classroomqa}
\lecturetimestamp{01:48:07}
\audiencequestion{最終散乱面が地平線へと薄れていくのは、銀河が地平線の向こうへ移動していくのと同じことなのか。}
\lecturerresponse{いいえ。後から来る放射は、今見ている放射と同じ原子から来ているのではない。それはもっと遠くにあった原子から来ている。したがって、放射がますます冷たく見える一方で、宇宙のますます深くを見ているようなものである。}
\end{classroomqa}

\section{曲率、ビッグクランチ像、そして後期の地平線}

ここまでは曲率の項を抑えてきた。それを復活させると
\begin{equation}
\left(\frac{\dot a}{a}\right)^2
=
\frac{8\pi G}{3}\rho
-
\frac{K}{a^2}.
\end{equation}
となる。数値的にはこの項は小さいように見えるが、概念的には重要である。古いニュートン的な言葉では、正、ゼロ、負の全エネルギーの場合を区別する。相対論的な言葉では、空間の幾何を区別する。

真空エネルギーが真剣に受け止められる以前は、膨張がいずれ逆転する物質優勢の宇宙がしばしば想像されていた。その結果が古いビッグクランチ像である。ビッグバンからの膨張、そして再収縮である。数学は今でも有効であるが、もはや我々の宇宙を記述しているようには見えない。

\begin{classroomqa}
\lecturetimestamp{01:53:51}
\audiencequestion{この真空エネルギーは量子真空エネルギーと関係があるのか。}
\lecturerresponse{はい、少なくともそう信じられている。}
\end{classroomqa}

現在のおおよその組成は
\begin{equation}
\frac{\rho_0}{\rho_{\text{total}}}\sim 0.7,
\qquad
\frac{\rho_{\text{matter}}}{\rho_{\text{total}}}\sim 0.3.
\end{equation}
である。その物質成分のうち、発光物質は全体の約\(0.05\)、暗黒物質は約\(0.25\)である。放射は現在無視できる。暗黒物質と発光物質は同じようにスケールするので、物質、放射、真空の相対的な存在量が時代ごとに変化しても、それらの比は固定されたままである。

真空優勢の極限では
\begin{equation}
H=\sqrt{\frac{8\pi G}{3}\rho_0}.
\end{equation}
となる。するとハッブルの法則は
\begin{equation}
v=Hd.
\end{equation}
となる。地平線は\(v=1\)で定義されるので、その距離は
\begin{equation}
d_H=\frac{1}{H}
=
\frac{1}{\sqrt{8\pi G\rho_0/3}}.
\end{equation}
である。ここで与えられたおおよその数は約\(10\)から\(12\)億光年である。

重要な点は後期の振る舞いである。膨張が指数関数的になると、地平線距離は固定される。銀河はそれを通過し続け、ますます赤方偏移して見えなくなっていく。宇宙マイクロ波背景放射もまた赤方偏移して見えなくなる。遠い未来には、近くの重力で束縛された構造だけが見え続け、宇宙は孤立した島のように見える。

\section{最終散乱面からの幾何}

最後に、幾何は最終散乱面から検証できる。その面までの距離が分かっているとしよう。また、そこに物理的な大きさが十分よく分かっている塊を同定できるとしよう。すると、その大きさと距離を観測された角度サイズと比較できる。

平坦な空間では角度はユークリッド的なものである。正に曲がった空間ではより大きく見え、負に曲がった空間ではより小さく見える。類推は普通の測量である。既知の距離にある既知の大きさの物体を取り、測定された角度をユークリッド幾何と比較する。

\begin{classroomqa}
\lecturetimestamp{02:08:09}
\audiencequestion{そこにある大きさが分かっている物体の例は何か。}
\lecturerresponse{宇宙マイクロ波背景放射の最大の塊である。我々はそれらが進化するのにかかった時間と音速から、その大きさを知っていると考えている。}
\end{classroomqa}

ここで関係する物体は、宇宙マイクロ波背景放射上の最大の塊である。それらの観測された角度サイズは、数パーセント以内でユークリッド幾何と一致する。それに基づくと、空間は平坦に見える。

\section{測定された\(w\)と最後の問い}

\(w\)の観測的な重要性は今や明らかである。天文学者が\(w\)は\(-1\)に非常に近いと報告するなら、彼らは真空エネルギー成分がほとんど一定であると言っているのである。それはまさにここで述べられた点である。

それは一般相対論が\(w\)を厳密に\(-1\)に強制するという意味ではない。ここでの点は経験的なものである。講義は、より急進的な可能性が原理的に排除されないことを簡単に述べるが、ほとんど一定の真空エネルギーの場合こそが信じるに足る良い理由がある場合であると述べている。

\begin{classroomqa}
\lecturetimestamp{02:11:04}
\audiencequestion{物質と放射の密度計算では、エネルギーが一定であると仮定してから、式を求めるためにスケール因子を調整したのですか。そして、それが最初の二つの場合に当てはまるなら、密度が一定に保たれる真空エネルギーの場合はどうなるのですか。}
\lecturerresponse{いいえ。私はそれが一定であるとは仮定しませんでした。我々はエネルギーの変化が\(-P\,dV\)であることを用いました。膨張とともに動く仮想的な共動箱に対して、壁で仕事がなされ、それがエネルギーの変化を説明します。真空エネルギーの場合はこの問いを別の形で提起し、ここでの結びの言葉は、宇宙論におけるエネルギー保存の言葉遣いはあまり多くを語らない、というものです。なぜなら宇宙の全エネルギーはゼロだからです。}
\end{classroomqa}
